\section{Implications for human play}
\label{app:human}

This project exists because a person plays this game and wanted a better opponent. This appendix
answers the questions that came from live play, in the order they were asked, and separates what
the evidence supports from what it does not.

\subsection{Does FishBot play for itself or for the team?}
\label{app:human-labels}

For the team, unambiguously, and this is a property of the objective rather than a tuning choice.
Every quantity the value function reads is team-relative: the score differential is the team's, the
card counts are the team's totals, and the per-half-suit control term is the expected fraction of
that half-suit held by the \emph{team}, not by the seat. There is no per-seat score in this game and
none in the policy. A configuration that maximised the acting seat's own holdings would be a
different program, and it is not the one that ships.

What is true, and is the substantive half of the question, is that the \emph{policy} is unilateral.
Through v0.5, FishBot maintained a deduction state for itself and for nobody else --- the only two
clones of that object anywhere in the policy are clones of its own --- so no term could ask what a
partner knows, or which of two teammates is better placed to act. Teammates entered only as
aggregate probability mass. v0.6 adds the machinery for per-seat modelling, uses it inside the
search, and reports that at the resolution of \S\ref{sec:protocol} the terms built on it are worth
nothing. That is a finding about this game rather than a concession: with every card movement
public, most of what a partner knows, you already know.

\subsection{Why a well-trained human still out-reasons it}
\label{app:human-ask}

The honest answer from this study is that the gap is in the \emph{action space}, not in the
arithmetic. M1, the mechanism that fixed v0.4's deadlock, restricts the candidate set to asks the
actor cannot prove will miss. That is the right fix for the deadlock and it deletes the move class
that most named human tactics are built from: blackballing, deliberately handing the turn to the
opponent who can do least with it, and paying for a signal with a safe ask are all a guaranteed miss
at a \emph{chosen} seat. Such a move is available at $79.2\%$ of decisions, and at two or more
distinct chosen opponents at $50.1\%$.

\S\ref{sec:search-deadask} reports what happened when it was given back. Unbanned wholesale, it
raises the win rate and destroys the policy. Rationed against the fitted linear score, it is never
chosen --- and the reason is structural rather than a matter of weights: the hit-probability term is
zero on a deliberate miss by construction, and every other term of the score is a penalty, so no
setting of the admission margin changes the output. Its value is the value of the position the
donated turn produces, which is a forward-looking quantity that no static feature computes. Offered
to the search, which is the only evaluator in the system that can price it, it produces the best
search configuration measured, and the search as a whole is the one mechanism in this study that
beats a static rule (\S\ref{sec:search}).

So: a human who blackballs is doing something the bot structurally cannot express, and the bot is
not obviously losing points by failing to. Both halves of that sentence are measured.

\subsection{Deception}
\label{app:human-bluff}

The manoeuvre the project owner reported --- holding cards of a half-suit you were asked for, then
declining to ask back in it --- is implemented as the \emph{withholder} archetype and is the one
opponent against which v0.6 is decisively better. The gain replicates in sign and size on two
disjoint seed banks and is the largest replicated effect in this study. It comes from the refit
rather than from an opponent model: the fitting panel contains the archetype and the objective is
minimax regret, so the optimiser buys robustness against it directly.

The complementary archetype, the \emph{feint}, manufactures a misleading ask-legality certificate
instead, and it remains the worst cell in the profile. v0.5's exposure there was created by its own
fit and could not be tuned away at fixed weights; v0.6's fit includes it in the panel and improves
it on one bank while giving a little back on the other.

\subsection{Playing against it}
\label{app:human-declare}

Three things a person should know when sitting down against this program.

\begin{itemize}
\item \textbf{It never asks a question it can prove will fail.} That is a hard deduction from your
  own public record, so a question it does ask is always one it cannot rule out --- which means an
  ask carries less information about its hand than a human's ask does.
\item \textbf{It is deterministic.} Given the same public history and the same hand it plays the
  same move. A person who has played many games against it can, in principle, read it. Adding
  stochasticity is the outstanding item: the machinery for a partner-aware temperature is specified
  and the harness now supports the two partner regimes as separate columns, but the temperature
  itself is not fitted here.
\item \textbf{Its declarations are conservative and accurate.} In mirror play it declares wrongly on
  roughly two percent of declarations, and the pre-gates that decide when to evaluate a declaration
  are exact --- \vsixGateFalseNeg{} false negatives in \vsixGateRejected{} rejections. If it
  declares, it is almost always right; the exploitable side is the timing, not the accuracy.
\end{itemize}

\subsection{Where a human should expect to keep an edge}
\label{app:human-ties}

At more than half of its decisions the program's own score cannot separate its leading candidates,
and \S\ref{sec:diag-irreducible} establishes that a correct posterior cannot separate them either.
A human facing the same information is in the same position. This is worth stating plainly because
it is the opposite of the usual intuition about card-game programs: the residual uncertainty here is
not a modelling deficiency that better software will remove --- it is the game.

Where a human should expect an edge is in the places this study found and could not close: choosing
\emph{which} opponent to hand the turn to when nothing good is available, paying a small material
price for a signal a partner will read, and varying behaviour so as not to be read. All three are
moves in the class M1 deletes.
